% Options for packages loaded elsewhere
\PassOptionsToPackage{unicode}{hyperref}
\PassOptionsToPackage{hyphens}{url}
\documentclass[
  letterpaper,
]{article}
\usepackage{xcolor}
\usepackage[margin=0.8in]{geometry}
\usepackage{amsmath,amssymb}
\setcounter{secnumdepth}{-\maxdimen} % remove section numbering
\usepackage{iftex}
\ifPDFTeX
  \usepackage[T1]{fontenc}
  \usepackage[utf8]{inputenc}
  \usepackage{textcomp} % provide euro and other symbols
\else % if luatex or xetex
  \usepackage{unicode-math} % this also loads fontspec
  \defaultfontfeatures{Scale=MatchLowercase}
  \defaultfontfeatures[\rmfamily]{Ligatures=TeX,Scale=1}
\fi
\usepackage{lmodern}
\ifPDFTeX\else
  % xetex/luatex font selection
\fi
% Use upquote if available, for straight quotes in verbatim environments
\IfFileExists{upquote.sty}{\usepackage{upquote}}{}
\IfFileExists{microtype.sty}{% use microtype if available
  \usepackage[]{microtype}
  \UseMicrotypeSet[protrusion]{basicmath} % disable protrusion for tt fonts
}{}
\makeatletter
\@ifundefined{KOMAClassName}{% if non-KOMA class
  \IfFileExists{parskip.sty}{%
    \usepackage{parskip}
  }{% else
    \setlength{\parindent}{0pt}
    \setlength{\parskip}{6pt plus 2pt minus 1pt}}
}{% if KOMA class
  \KOMAoptions{parskip=half}}
\makeatother
\usepackage{listings}
\newcommand{\passthrough}[1]{#1}
\lstset{defaultdialect=[5.3]Lua}
\lstset{defaultdialect=[x86masm]Assembler}
\setlength{\emergencystretch}{3em} % prevent overfull lines
\providecommand{\tightlist}{%
  \setlength{\itemsep}{0pt}\setlength{\parskip}{0pt}}
% Preserve readable code, command, and table layout in Pandoc-generated review PDFs.
\usepackage{etoolbox}
\usepackage{pdflscape}
\lstset{
  basicstyle=\ttfamily\footnotesize,
  breaklines=true,
  breakatwhitespace=false,
  columns=fullflexible,
  keepspaces=true,
  showstringspaces=false,
  upquote=true
}
\AtBeginEnvironment{longtable}{\footnotesize}
\renewcommand{\arraystretch}{1.08}
\setlength{\emergencystretch}{3em}
\usepackage{bookmark}
\IfFileExists{xurl.sty}{\usepackage{xurl}}{} % add URL line breaks if available
\urlstyle{same}
% fallback for those not using the hyperref driver hyperxmp:
\makeatletter
\@ifundefined{xmpquote}{\newcommand{\xmpquote}[1]{#1}}{}
\makeatother
\hypersetup{
  pdftitle={Password Bonobo URL Audit and Cleanup --- Design Specification},
  hidelinks,
  pdfcreator={LaTeX via pandoc}}

\title{Password Bonobo URL Audit and Cleanup --- Design Specification}
\author{}
\date{}

\begin{document}
\maketitle

{
\setcounter{tocdepth}{3}
\tableofcontents
}
\section{Password Bonobo URL Audit and Cleanup --- Design
Specification}\label{password-bonobo-url-audit-and-cleanup--design-specification}

Date: 2026-08-22 Status: Approved design, pending implementation plan Target: Password Bonobo (native Tcl/Tk
implementation based on Password Gorilla)

\subsection{1. Purpose}\label{1-purpose}

Add a native Password Bonobo feature, based on Password Gorilla, that audits login entries with populated URL fields,
identifies websites that are likely obsolete or require human review, and supports safe bulk cleanup of obsolete
credentials.

The primary user goal is to reduce the size of a large Password Gorilla database by finding credentials for websites
that are no longer useful and deleting them deliberately. The feature must prioritize reviewability and data safety over
aggressive automatic classification.

\subsection{2. Core Principles}\label{2-core-principles}

\begin{enumerate}
\def\labelenumi{\arabic{enumi}.}
\tightlist
\item
  The scanner never automatically deletes or modifies password entries.
\item
  Network classification is advisory. The user makes the final decision.
\item
  Ambiguous results are classified as \passthrough{\lstinline!Needs Review!}, not \passthrough{\lstinline!Dead!}.
\item
  Deletion changes only the in-memory database until the user invokes Gorilla\textquotesingle s normal Save command.
\item
  Archiving is an explicit export to a separate encrypted PasswordSafe/Gorilla database file, not an ``archived'' flag
  on an entry.
\item
  \passthrough{\lstinline!Archive \& Delete Selected!} is the primary cleanup action;
  \passthrough{\lstinline!Delete Selected!} is secondary.
\item
  Archive creation must succeed completely before any selected entry is deleted from memory.
\item
  The scanner must not transmit usernames, passwords, notes, query strings, fragments, or other credential contents.
\end{enumerate}

\subsection{3. Scope}\label{3-scope}

\subsubsection{3.1 In scope}\label{31-in-scope}

\begin{itemize}
\tightlist
\item
  Scan every login entry with a nonempty URL field.
\item
  Normalize web URLs for testing without modifying stored values.
\item
  Strip query strings and fragments from the network request.
\item
  Check URLs asynchronously while keeping the Tk UI responsive.
\item
  Follow a bounded number of redirects.
\item
  Classify results as \passthrough{\lstinline!Working!}, \passthrough{\lstinline!Needs Review!},
  \passthrough{\lstinline!Dead!}, or \passthrough{\lstinline!Not Checked!}.
\item
  Detect common parked/domain-for-sale pages conservatively.
\item
  Show URL hygiene warnings independently from website-health status.
\item
  Present results in a modeless review window.
\item
  Allow browser opening, jumping to the corresponding Gorilla entry, and rechecking individual rows.
\item
  Support bulk selection and deletion.
\item
  Support exporting selected records to a new encrypted \passthrough{\lstinline!.psafe3!} archive before deletion.
\item
  Use the current database master password by default for the archive, with an option to choose a different password.
\item
  Preserve the original on-disk database until the user explicitly saves it.
\end{itemize}

\subsubsection{3.2 Out of scope for version 1}\label{32-out-of-scope-for-version-1}

\begin{itemize}
\tightlist
\item
  Automatically deleting entries.
\item
  Automatically saving the cleaned main database.
\item
  Clearing only the URL field.
\item
  Adding an archive flag to the PasswordSafe record format.
\item
  Appending to an existing archive database.
\item
  Overwriting an existing archive database.
\item
  Automatically rewriting stored URLs after redirects.
\item
  Logging in to websites or submitting credentials.
\item
  Executing JavaScript or behaving as a full browser.
\item
  Following non-web schemes such as \passthrough{\lstinline!ftp:!}, \passthrough{\lstinline!ssh:!}, or custom
  application schemes.
\end{itemize}

\subsection{4. Entry Point}\label{4-entry-point}

Add a command under the existing Login menu:

\passthrough{\lstinline!Login -> Audit Website URLs...!}

The command is available only when a database is open.

A new top-level Tools menu is not introduced for version 1.

\subsection{5. URL Collection and Normalization}\label{5-url-collection-and-normalization}

For each login entry:

\begin{enumerate}
\def\labelenumi{\arabic{enumi}.}
\tightlist
\item
  Read the URL field.
\item
  Skip entries whose URL is empty or whitespace only.
\item
  Preserve the original URL exactly as stored for display and later browser opening.
\item
  Construct a separate audit URL used only for network testing.
\item
  Remove the fragment (\passthrough{\lstinline!\#...!}) from the audit URL.
\item
  Remove the query string (\passthrough{\lstinline!?...!}) from the audit URL.
\item
  If the URL has an explicit \passthrough{\lstinline!https://!} scheme, test it as written after sanitization.
\item
  If the URL has an explicit \passthrough{\lstinline!http://!} scheme, test it as written after sanitization and permit
  a normal redirect to HTTPS.
\item
  If no scheme is present, try \passthrough{\lstinline!https://!} first. Only if that fails in a way consistent with
  there being no usable HTTPS endpoint may the checker attempt \passthrough{\lstinline!http://!}.
\item
  If an explicit non-HTTP(S) scheme is present, classify the entry \passthrough{\lstinline!Not Checked!}.
\item
  Malformed values that cannot be safely normalized are \passthrough{\lstinline!Not Checked!}, never
  \passthrough{\lstinline!Dead!}.
\end{enumerate}

Normalization never changes the URL stored in the database.

\subsection{6. URL Hygiene Warnings}\label{6-url-hygiene-warnings}

Website health and URL hygiene are separate dimensions.

Possible hygiene notes include:

\begin{itemize}
\tightlist
\item
  \passthrough{\lstinline!Query parameters!}
\item
  \passthrough{\lstinline!Fragment!}
\item
  \passthrough{\lstinline!Potentially sensitive parameters!}
\end{itemize}

A URL remains \passthrough{\lstinline!Working!} if the website works even when a hygiene warning is present.

Potentially sensitive query parameter names should be detected case-insensitively using a conservative list such as:

\begin{itemize}
\tightlist
\item
  \passthrough{\lstinline!token!}
\item
  \passthrough{\lstinline!auth!}
\item
  \passthrough{\lstinline!key!}
\item
  \passthrough{\lstinline!session!}
\item
  \passthrough{\lstinline!code!}
\item
  \passthrough{\lstinline!password!}
\item
  \passthrough{\lstinline!passwd!}
\item
  \passthrough{\lstinline!reset!}
\item
  \passthrough{\lstinline!email!}
\item
  \passthrough{\lstinline!user!}
\item
  \passthrough{\lstinline!username!}
\item
  \passthrough{\lstinline!account!}
\end{itemize}

Only parameter names are inspected for warning purposes. Parameter values must not be logged, transmitted, or copied
into diagnostic output.

\subsection{7. Network Checking Model}\label{7-network-checking-model}

\subsubsection{7.1 Native Tcl implementation}\label{71-native-tcl-implementation}

Implement the checker inside Gorilla using Tcl\textquotesingle s asynchronous HTTP facilities rather than invoking an
external executable such as \passthrough{\lstinline!curl!}.

The implementation should use a small bounded pool of concurrent requests. A default concurrency of 4 is recommended;
the internal design may permit a small maximum such as 6 if testing shows that this remains responsive and polite.

\subsubsection{7.2 HTTPS requirement}\label{72-https-requirement}

HTTPS checking requires a TLS-capable Tcl socket provider. Password Gorilla\textquotesingle s source bundle includes
Tcllib, but TclTLS is a separate extension. The implementation must therefore:

\begin{itemize}
\tightlist
\item
  require and configure TclTLS (or an equivalent native Tcl TLS provider);
\item
  preserve certificate validation;
\item
  support SNI;
\item
  ship or otherwise reliably provide the appropriate CA trust material on supported platforms;
\item
  fail closed if HTTPS cannot be validated rather than silently disabling certificate verification.
\end{itemize}

If TLS support is unavailable, the audit should display a clear prerequisite/error message rather than misclassifying
HTTPS sites as dead.

\subsubsection{7.3 Request behavior}\label{73-request-behavior}

\begin{itemize}
\tightlist
\item
  Use ordinary unauthenticated HTTP(S) GET requests.
\item
  Do not send cookies, usernames, passwords, notes, or database metadata.
\item
  Use a neutral, non-identifying User-Agent that does not reveal the database name, entry title, username, or
  master-password-manager context.
\item
  Do not rely solely on HEAD because many servers handle HEAD differently from GET and parked-page detection requires
  examining a small response body.
\item
  Stop reading response bodies once enough data has been obtained for classification; do not download large pages
  unnecessarily.
\item
  Apply a bounded per-request timeout. The initial implementation should use a value around 10 seconds and make the
  constant easy to tune.
\item
  Follow at most 5 redirects.
\item
  Detect and terminate redirect loops.
\end{itemize}

\subsection{8. Classification Rules}\label{8-classification-rules}

The scanner exposes four health classifications.

\subsubsection{8.1 Working}\label{81-working}

Classify as \passthrough{\lstinline!Working!} when there is strong evidence that the saved site is still usable,
including:

\begin{itemize}
\tightlist
\item
  successful 2xx response;
\item
  ordinary redirect that resolves successfully;
\item
  HTTP-to-HTTPS upgrade;
\item
  redirect within the same hostname;
\item
  conservative canonical-host normalization such as adding or removing only a leading \passthrough{\lstinline!www.!}.
\end{itemize}

A redirect to a meaningfully different hostname is not automatically considered working even if the final response is
successful.

\subsubsection{8.2 Needs Review}\label{82-needs-review}

Use \passthrough{\lstinline!Needs Review!} whenever the evidence is ambiguous. Examples include:

\begin{itemize}
\tightlist
\item
  \passthrough{\lstinline!401 Unauthorized!};
\item
  \passthrough{\lstinline!403 Forbidden!};
\item
  \passthrough{\lstinline!429 Too Many Requests!};
\item
  CAPTCHA, bot protection, or Cloudflare-style challenge;
\item
  TLS/certificate errors;
\item
  timeout;
\item
  connection refusal;
\item
  \passthrough{\lstinline!5xx!} server error;
\item
  redirect to a substantially different hostname;
\item
  saved path returns \passthrough{\lstinline!404!} or \passthrough{\lstinline!410!} while the site root remains alive;
\item
  page content suggests possible parking/retirement but does not meet the confidence threshold for
  \passthrough{\lstinline!Dead!};
\item
  any result where the checker cannot distinguish a dead service from a temporary or automation-specific failure.
\end{itemize}

A saved login page returning \passthrough{\lstinline!404!} or \passthrough{\lstinline!410!} does not by itself make the
entry \passthrough{\lstinline!Dead!}. The checker must test the site root before deciding whether the site as a whole
appears gone.

\subsubsection{8.3 Dead}\label{83-dead}

\passthrough{\lstinline!Dead!} is intentionally high-confidence. It may be assigned when, for example:

\begin{itemize}
\tightlist
\item
  DNS resolution definitively reports that the hostname does not exist;
\item
  the target or redirect destination matches a confidently recognized domain-parking/domain-for-sale pattern;
\item
  both the saved target and the site root return evidence that the site has intentionally been retired, such as a
  conclusive \passthrough{\lstinline!410 Gone!}, with no contrary evidence.
\end{itemize}

Temporary network failures must not be promoted to \passthrough{\lstinline!Dead!} merely because retries failed.

\subsubsection{8.4 Not Checked}\label{84-not-checked}

Use \passthrough{\lstinline!Not Checked!} for:

\begin{itemize}
\tightlist
\item
  non-HTTP(S) schemes;
\item
  malformed URLs that cannot be safely normalized;
\item
  unsupported URL forms;
\item
  other entries intentionally excluded from web checking.
\end{itemize}

\subsection{9. Parked / Domain-for-Sale Detection}\label{9-parked--domain-for-sale-detection}

Parking detection should be conservative and heuristic-driven.

The detector may inspect a bounded amount of returned HTML/text for recognizable patterns associated with registrar
parking, domain-for-sale pages, or known parking providers.

Requirements:

\begin{itemize}
\tightlist
\item
  Keep parking rules isolated from general HTTP classification so they can be expanded or corrected without changing the
  scanner core.
\item
  Prefer false negatives over false positives.
\item
  If confidence is not high, classify \passthrough{\lstinline!Needs Review!} rather than \passthrough{\lstinline!Dead!}.
\item
  Never make deletion automatic based on parking detection.
\end{itemize}

\subsection{10. Review Window}\label{10-review-window}

Open a modeless \passthrough{\lstinline!URL Audit Results!} window so the user can continue interacting with the main
Gorilla window.

\subsubsection{10.1 Progress area}\label{101-progress-area}

Show:

\begin{itemize}
\tightlist
\item
  \passthrough{\lstinline!Scanning: N / Total!}
\item
  progress bar
\item
  \passthrough{\lstinline!Cancel Scan!} button while scanning
\end{itemize}

Results appear incrementally as checks complete.

The user may begin reviewing completed results before the full scan finishes.

Cancellation stops launching new checks and cancels outstanding checks where possible. Completed results remain
available.

\subsubsection{10.2 Result table}\label{102-result-table}

Use a Tk/ttk multi-column list/tree control with columns equivalent to:

\begin{itemize}
\tightlist
\item
  selected checkbox/state
\item
  Group / Entry
\item
  Stored URL
\item
  Status
\item
  Reason
\item
  Redirect / URL warning
\end{itemize}

The model should retain, per row:

\begin{itemize}
\tightlist
\item
  record UUID or stable record identifier;
\item
  group;
\item
  title;
\item
  original stored URL;
\item
  sanitized audit URL;
\item
  status;
\item
  detailed reason;
\item
  final redirect destination, if any;
\item
  hygiene warning(s);
\item
  deletion/archive state.
\end{itemize}

\subsubsection{10.3 Default filtering}\label{103-default-filtering}

Show by default:

\begin{itemize}
\tightlist
\item
  \passthrough{\lstinline!Dead!}
\item
  \passthrough{\lstinline!Needs Review!}
\end{itemize}

Hide by default:

\begin{itemize}
\tightlist
\item
  \passthrough{\lstinline!Working!}
\item
  \passthrough{\lstinline!Not Checked!}
\end{itemize}

Provide filter controls for all four categories.

\subsubsection{10.4 Selection behavior}\label{104-selection-behavior}

\begin{itemize}
\tightlist
\item
  Nothing is selected for deletion automatically.
\item
  Provide \passthrough{\lstinline!Select All Dead!}.
\item
  Provide \passthrough{\lstinline!Select Visible!} if straightforward in the chosen widget implementation.
\item
  Provide \passthrough{\lstinline!Clear Selection!}.
\item
  Selection is independent from classification; the user may choose any visible result.
\end{itemize}

\subsubsection{10.5 Row actions}\label{105-row-actions}

Provide row-level actions:

\begin{itemize}
\tightlist
\item
  \passthrough{\lstinline!Open Website!} --- opens the original stored URL in the user\textquotesingle s
  configured/default browser using Gorilla\textquotesingle s existing browser-launch behavior;
\item
  \passthrough{\lstinline!Go to Login!} --- selects/reveals the corresponding entry in Gorilla\textquotesingle s main
  tree;
\item
  \passthrough{\lstinline!Recheck!} --- repeats the network classification for that record.
\end{itemize}

Double-clicking a row should open the stored website in the browser unless platform/UI conventions strongly favor
another existing Gorilla behavior.

\subsubsection{10.6 Bottom actions}\label{106-bottom-actions}

Primary action:

\passthrough{\lstinline!Archive \& Delete Selected!}

Secondary action:

\passthrough{\lstinline!Delete Selected!}

Other action:

\passthrough{\lstinline!Close!}

After deletion, keep affected rows in the audit window and mark them \passthrough{\lstinline!Deleted (unsaved)!} rather
than removing them immediately. This preserves a visible audit trail for the current session.

\subsection{11. Deleting Selected Entries}\label{11-deleting-selected-entries}

\passthrough{\lstinline!Delete Selected!} performs the following:

\begin{enumerate}
\def\labelenumi{\arabic{enumi}.}
\tightlist
\item
  Validate the selected records still exist in the current in-memory database.
\item
  Detect protected entries and exclude them from deletion.
\item
  Ask for explicit confirmation, including the count of entries to be deleted.
\item
  Delete confirmed, unprotected entries from the in-memory database only.
\item
  Mark the current database modified using Gorilla\textquotesingle s normal modified-state mechanism.
\item
  Update the audit rows to \passthrough{\lstinline!Deleted (unsaved)!}.
\item
  Do not invoke Save.
\end{enumerate}

The existing on-disk database remains unchanged until the user uses Gorilla\textquotesingle s normal Save command.

If the user closes Gorilla without saving and chooses not to save changes, the original database on disk remains intact.

\subsection{12. Archive \& Delete Selected}\label{12-archive--delete-selected}

Archiving is implemented as creation of a separate encrypted PasswordSafe/Gorilla database file containing the selected
records.

\subsubsection{12.1 Archive file behavior}\label{121-archive-file-behavior}

\begin{itemize}
\item
  Create a brand-new \passthrough{\lstinline!.psafe3!} file.
\item
  Do not append to an existing archive in version 1.
\item
  Do not silently overwrite an existing file.
\item
  Default filename pattern:

  \nolinkurl{<current-database-base>-dead-YYYY-MM-DD.psafe3}
\item
  Allow the user to choose another filename/location.
\end{itemize}

\subsubsection{12.2 Archive password}\label{122-archive-password}

Default option:

\passthrough{\lstinline!Use current database master password!}

Alternative:

\passthrough{\lstinline!Use a different archive password!}

If a different password is selected, require it to be entered twice and validate that the two values match before
writing.

The URL-audit feature must not cause Gorilla to retain the current master password longer than it already does. If
Gorilla does not retain the current master password in retrievable form after unlock, choosing
\passthrough{\lstinline!Use current database master password!} should prompt the user to re-enter that password rather
than introducing new long-lived password storage.

\subsubsection{12.3 Transactional sequence}\label{123-transactional-sequence}

\passthrough{\lstinline!Archive \& Delete Selected!} must perform these steps in order:

\begin{enumerate}
\def\labelenumi{\arabic{enumi}.}
\tightlist
\item
  Snapshot/copy the complete selected records required for archive creation.
\item
  Re-resolve the selected records by stable identifier and validate protected-entry behavior before destructive changes.
  If protected or stale records are present, present an explicit summary of which records will be skipped before
  proceeding. Protected records are neither archived nor deleted unless the user first unprotects them through
  Gorilla\textquotesingle s normal UI.
\item
  Create a new PasswordSafe database object for the archive.
\item
  Copy the complete selected records into it, preserving record fields and stable record identity where the existing
  PasswordSafe implementation permits this safely.
\item
  Write the archive to a temporary file in the destination directory.
\item
  Close the archive successfully.
\item
  Reopen/validate the archive using the selected archive password and verify that the expected record count and record
  identities are present.
\item
  Atomically rename the temporary file to the requested archive filename.
\item
  Only after steps 1--8 succeed, delete the selected unprotected entries from the current in-memory database.
\item
  Mark the current database modified.
\item
  Mark result rows \passthrough{\lstinline!Archived \& deleted (unsaved)!}.
\item
  Do not save the current main database automatically.
\end{enumerate}

If any archive step fails, no selected entry is deleted.

\subsubsection{12.4 Protected entries}\label{124-protected-entries}

Entries protected by PasswordSafe/Gorilla protections may appear in audit results but must not be silently unprotected
or deleted by the audit tool.

If selected, the operation should identify them and tell the user they must first be unprotected through
Gorilla\textquotesingle s normal entry-editing mechanism. A mixed selection may proceed only after Gorilla clearly
reports the count of protected/stale records being skipped and the count that will actually be archived/deleted.

\subsection{13. Data Safety and Failure Handling}\label{13-data-safety-and-failure-handling}

\subsubsection{13.1 Scan failures}\label{131-scan-failures}

A network failure never changes the database.

\subsubsection{13.2 Archive failures}\label{132-archive-failures}

Examples include:

\begin{itemize}
\tightlist
\item
  invalid destination;
\item
  permission denied;
\item
  write failure;
\item
  disk-full condition;
\item
  encryption/database-write error;
\item
  archive validation failure;
\item
  rename failure.
\end{itemize}

Any such failure leaves the current database entries untouched.

Temporary archive files should be removed when safe to do so after a failed operation. A failure to remove a temporary
file should be reported without deleting original records.

\subsubsection{13.3 Stale results}\label{133-stale-results}

Because the review window is modeless, a record may be edited or deleted in the main Gorilla window after scanning.
Before any destructive action, re-resolve each selected record by stable identifier and validate that it still exists.

If a record no longer exists, mark it stale/skipped rather than treating that as a fatal error for unrelated selected
records.

\subsection{14. Security and Privacy}\label{14-security-and-privacy}

The audit window is security-sensitive because it exposes entry titles, groups, and URLs. It is bound to the database
session that created it. If the database is locked, closed, or replaced with another database, Gorilla must cancel
outstanding audit requests and close the audit window (or equivalently remove all sensitive row data and disable
actions). The preferred version-1 behavior is to close the audit window. Background scanning must not defeat
Gorilla\textquotesingle s normal idle-lock policy.

\begin{itemize}
\tightlist
\item
  Never transmit passwords, usernames, notes, or other credential fields.
\item
  Strip query strings and fragments before network requests.
\item
  Do not log sensitive URL parameter values.
\item
  Avoid logging full URLs or URL paths in routine diagnostics. If diagnostic logging is needed, prefer hostname plus
  non-sensitive status/error metadata.
\item
  Do not include the database filename, entry title, group, username, or UUID in HTTP headers.
\item
  Keep TLS certificate validation enabled.
\item
  Treat TLS validation failures as \passthrough{\lstinline!Needs Review!}.
\item
  Do not silently downgrade an explicitly \passthrough{\lstinline!https://!} stored URL to HTTP.
\item
  For a URL without a scheme, HTTPS is attempted first; HTTP fallback is only a discovery mechanism for scheme-less
  input.
\item
  Archive files are always encrypted PasswordSafe/Gorilla databases, never plaintext CSV/JSON/text exports.
\item
  Archive passwords must be handled with the same sensitivity as the current master password and cleared from temporary
  UI variables where Gorilla\textquotesingle s conventions permit.
\end{itemize}

\subsection{15. Proposed Internal Components}\label{15-proposed-internal-components}

Keep the implementation separated into small Tcl namespaces/modules rather than adding all logic directly to
\passthrough{\lstinline!gorilla.tcl!}.

Suggested responsibilities:

\subsubsection{\texorpdfstring{\texttt{gorilla::URLAudit}}{gorilla::URLAudit}}\label{url-audit}

Owns scan orchestration and audit state:

\begin{itemize}
\tightlist
\item
  collect eligible records;
\item
  schedule requests;
\item
  maintain bounded concurrency;
\item
  cancellation;
\item
  aggregate progress;
\item
  connect checker results to UI rows.
\end{itemize}

\subsubsection{\texorpdfstring{\texttt{gorilla::URLAudit::URL}}{gorilla::URLAudit::URL}}\label{url}

Pure URL-handling functions:

\begin{itemize}
\tightlist
\item
  sanitize query/fragment;
\item
  detect/normalize scheme;
\item
  classify non-web schemes;
\item
  derive root URL;
\item
  compare redirect hosts conservatively;
\item
  detect hygiene warnings.
\end{itemize}

\subsubsection{\texorpdfstring{\texttt{gorilla::URLAudit::HTTP}}{gorilla::URLAudit::HTTP}}\label{http}

Network boundary:

\begin{itemize}
\tightlist
\item
  asynchronous GET;
\item
  timeout;
\item
  redirect handling;
\item
  response-size cap;
\item
  TLS setup;
\item
  response normalization;
\item
  cleanup/cancellation.
\end{itemize}

\subsubsection{\texorpdfstring{\texttt{gorilla::URLAudit::Classifier}}{gorilla::URLAudit::Classifier}}\label{classifier}

Pure classification logic:

\begin{itemize}
\tightlist
\item
  map network observations to \passthrough{\lstinline!Working!}, \passthrough{\lstinline!Needs Review!},
  \passthrough{\lstinline!Dead!}, \passthrough{\lstinline!Not Checked!};
\item
  root fallback checks;
\item
  parking detection;
\item
  human-readable reason strings.
\end{itemize}

\subsubsection{\texorpdfstring{\texttt{gorilla::URLAudit::Dialog}}{gorilla::URLAudit::Dialog}}\label{dialog}

Review UI:

\begin{itemize}
\tightlist
\item
  modeless window;
\item
  progress controls;
\item
  result list;
\item
  filters;
\item
  selection controls;
\item
  row actions;
\item
  delete/archive buttons;
\item
  status updates.
\end{itemize}

\subsubsection{\texorpdfstring{\texttt{gorilla::URLAudit::Archive}}{gorilla::URLAudit::Archive}}\label{archive}

Archive transaction:

\begin{itemize}
\tightlist
\item
  create archive DB;
\item
  copy records;
\item
  temporary-file write;
\item
  validation reopen;
\item
  atomic rename;
\item
  only then request in-memory deletion.
\end{itemize}

Where practical, the modules should reuse existing Gorilla record-selection, browser-launch, database-modified,
PasswordSafe-write, and delete-entry functions rather than duplicating them.

\subsection{16. Testing Strategy}\label{16-testing-strategy}

\subsubsection{16.1 URL parsing / hygiene unit tests}\label{161-url-parsing--hygiene-unit-tests}

Cover:

\begin{itemize}
\tightlist
\item
  \passthrough{\lstinline!https://example.com/login!}
\item
  \passthrough{\lstinline!http://example.com/login!}
\item
  \passthrough{\lstinline!example.com/login!}
\item
  query strings;
\item
  fragments;
\item
  query plus fragment;
\item
  potentially sensitive parameter names;
\item
  malformed strings;
\item
  explicit non-web schemes;
\item
  Unicode/escaped URL edge cases where supported by the existing URL parser.
\end{itemize}

Verify the original stored URL is never modified.

\subsubsection{16.2 HTTP/classification tests}\label{162-httpclassification-tests}

Use a controlled local HTTP test server where possible. Cover:

\begin{itemize}
\tightlist
\item
  200;
\item
  204;
\item
  same-host redirect;
\item
  HTTP-to-HTTPS redirect where test infrastructure supports it;
\item
  \passthrough{\lstinline!www!} canonical redirect;
\item
  cross-host redirect;
\item
  401;
\item
  403;
\item
  404 with live root;
\item
  410 with live root;
\item
  conclusive retired root;
\item
  429;
\item
  500/502/503;
\item
  timeout;
\item
  connection refusal;
\item
  DNS failure through an injectable resolver/network boundary where deterministic testing is required;
\item
  redirect loop;
\item
  redirect limit exceeded;
\item
  TLS certificate failure;
\item
  parked-domain signatures;
\item
  uncertain parking-like page.
\end{itemize}

\subsubsection{16.3 UI tests}\label{163-ui-tests}

Verify:

\begin{itemize}
\tightlist
\item
  scan does not freeze the main Tk event loop;
\item
  progress count is correct;
\item
  results appear incrementally;
\item
  cancellation works;
\item
  completed results survive cancellation;
\item
  default filters show only \passthrough{\lstinline!Dead!} and \passthrough{\lstinline!Needs Review!};
\item
  nothing is selected automatically;
\item
  \passthrough{\lstinline!Select All Dead!}, \passthrough{\lstinline!Select Visible!}, and
  \passthrough{\lstinline!Clear Selection!} behave correctly;
\item
  browser opening uses the original stored URL;
\item
  \passthrough{\lstinline!Go to Login!} selects the correct main-tree entry;
\item
  recheck updates one row safely;
\item
  deleted rows remain visible as \passthrough{\lstinline!Deleted (unsaved)!}.
\end{itemize}

\subsubsection{16.4 Database safety tests}\label{164-database-safety-tests}

Verify:

\begin{itemize}
\tightlist
\item
  \passthrough{\lstinline!Delete Selected!} changes only the in-memory database;
\item
  the database is marked modified after deletion;
\item
  the original file bytes do not change before Save;
\item
  choosing not to save on exit leaves the original file unchanged;
\item
  protected entries are not deleted;
\item
  stale/deleted records are skipped safely.
\end{itemize}

\subsubsection{16.5 Archive transaction tests}\label{165-archive-transaction-tests}

Verify:

\begin{itemize}
\tightlist
\item
  archive contains exactly the selected records;
\item
  expected fields survive intact, including group, title, username, password, URL, notes, and supported metadata;
\item
  archive opens successfully with the chosen password;
\item
  archive can be opened by Gorilla and, where practical, Password Safe;
\item
  same-current-password path works;
\item
  different-password path works;
\item
  password confirmation mismatch prevents writing;
\item
  destination collision is handled safely;
\item
  permission failure deletes nothing;
\item
  simulated write failure deletes nothing;
\item
  simulated validation failure deletes nothing;
\item
  simulated rename failure deletes nothing;
\item
  successful archive is finalized before main-memory deletion occurs.
\end{itemize}

\subsubsection{16.6 Large-database test}\label{166-large-database-test}

Test against a synthetic database with enough URL-bearing records to expose UI freezes, runaway memory usage, or poor
scheduling behavior.

Success criteria:

\begin{itemize}
\tightlist
\item
  UI remains responsive throughout the scan;
\item
  concurrency never exceeds the configured bound;
\item
  cancellation is prompt;
\item
  memory growth remains proportional to result count rather than response-body size;
\item
  review and bulk selection remain usable with hundreds or thousands of rows.
\end{itemize}

\subsection{17. Acceptance Criteria}\label{17-acceptance-criteria}

The feature is complete when all of the following are true:

\begin{enumerate}
\def\labelenumi{\arabic{enumi}.}
\tightlist
\item
  A user can invoke \passthrough{\lstinline!Login -> Audit Website URLs...!} on an open database.
\item
  Every nonempty URL field is either checked or explicitly classified \passthrough{\lstinline!Not Checked!}.
\item
  Query strings and fragments are never sent in audit HTTP requests.
\item
  The main UI remains responsive while scanning.
\item
  Results appear incrementally and the scan can be cancelled.
\item
  Dead-site classification remains conservative; ambiguous cases become \passthrough{\lstinline!Needs Review!}.
\item
  Working entries are hidden by default but can be shown.
\item
  No entry is preselected for deletion.
\item
  The user can open the stored URL, jump to the login entry, and recheck a result.
\item
  \passthrough{\lstinline!Delete Selected!} affects memory only and never saves automatically.
\item
  \passthrough{\lstinline!Archive \& Delete Selected!} creates and validates a new encrypted archive before deleting
  from memory.
\item
  Archive failure leaves all selected entries untouched.
\item
  Protected entries are respected.
\item
  The original on-disk database changes only when the user explicitly saves it through Gorilla\textquotesingle s normal
  Save workflow.
\item
  Generated archives open successfully as PasswordSafe/Gorilla databases.
\end{enumerate}

\subsection{18. Repository / Format Context}\label{18-repository--format-context}

The design is aligned with the current public Password Gorilla source structure:

\begin{itemize}
\tightlist
\item
  Gorilla is a Tcl/Tk application whose main UI and menus are in \passthrough{\lstinline!sources/gorilla.tcl!}.
\item
  The source bundle includes Tcllib and a PasswordSafe implementation under \passthrough{\lstinline!sources/pwsafe!}.
\item
  Gorilla\textquotesingle s existing UI already uses \passthrough{\lstinline!ttk::treeview!} and has normal Login-menu
  delete behavior, browser-launch behavior, and File Save/Save As/Export commands that the feature should reuse rather
  than duplicate.
\item
  PasswordSafe V3 is an encrypted, extensible, platform-independent format with per-record fields including UUID and
  URL-related record data.
\item
  TclTLS supports Tcl 8.5+ but requires explicit TLS/certificate configuration; HTTPS support must therefore be treated
  as a packaging/security requirement rather than assumed. Because current Gorilla source includes Tcl 9 compatibility
  work, packaged TclTLS builds must also match the Tcl major version used by each distribution.
\end{itemize}

References:

\begin{itemize}
\tightlist
\item
  \url{https://github.com/zdia/gorilla}
\item
  \url{https://github.com/zdia/gorilla/blob/master/sources/gorilla.tcl}
\item
  \url{https://github.com/pwsafe/pwsafe/blob/master/docs/formatV3.txt}
\item
  \url{https://core.tcl-lang.org/tcltls/doc/trunk/doc/tls.html}
\end{itemize}

\subsection{19. Explicitly Deferred Enhancements}\label{19-explicitly-deferred-enhancements}

Potential later enhancements, not required for version 1:

\begin{itemize}
\tightlist
\item
  append selected entries to an existing archive;
\item
  persistent audit history;
\item
  user-configurable timeout/concurrency settings;
\item
  richer domain-parking provider database;
\item
  optional URL-update suggestions after redirects;
\item
  re-run only previously ambiguous/dead results;
\item
  export an audit report that contains no secrets;
\item
  scheduled/automatic audits;
\item
  more sophisticated registrable-domain comparison using a Public Suffix List;
\item
  user-maintained allow/ignore rules for known sites.
\end{itemize}

\end{document}
